\chapter*{Introduction générale}
\addcontentsline{toc}{chapter}{Introduction générale}
\markboth{Introduction générale}{}

Le vent conditionne une grande partie des activités humaines exposées au climat : la réussite d'un
traitement phytosanitaire, la quantité d'eau perdue par une parcelle irriguée, la production d'une
éolienne, la sécurité d'un levage sur un chantier. Ces décisions se prennent à des hauteurs très
différentes --- 2 à 10\,m pour l'agriculture, 50 à 150\,m pour l'éolien --- et en des lieux où, le
plus souvent, aucun instrument ne mesure le vent. Les stations météorologiques sont rares, placées
près du sol, et les modèles numériques de prévision décrivent l'atmosphère sur des mailles de
plusieurs kilomètres qui lissent l'effet du relief et du littoral.

Tukhnanutha, holding technologique tunisienne, développe l'écosystème ANEO : un réseau de capteurs
environnementaux portés par des mâts multiservices, et des services numériques construits sur leurs
données. Ce projet de fin d'études répond à une question centrale de cet écosystème :
\emph{comment connaître le vent partout, à toutes les hauteurs utiles, à partir de mesures prises
en quelques points près du sol ?}

Nous y répondons par un modèle d'apprentissage profond contraint par la physique,
\textbf{PC-FNO3D}, un opérateur neuronal de Fourier tridimensionnel qui reconstruit en moins d'une
seconde les trois composantes du vent sur douze niveaux, de 25 à 1\,000\,m d'altitude. Le modèle
apprend sur des données dont il respecte la physique : la conservation de la masse et la condition
au sol font partie de sa fonction de perte. Il est ensuite transposé à un domaine réel, le Grand
Tunis, confronté aux mesures réelles des anémomètres et des radiosondages, et mis en service dans
une application web, \textbf{ATOM Wind}, destinée aux agriculteurs et aux exploitants éoliens.

Le travail s'est déroulé en plusieurs étapes qui structurent ce rapport :
\begin{itemize}
  \item le \textbf{chapitre~\ref{chap:contexte}} présente l'organisme d'accueil, le problème, les
        objectifs mesurables (DSO) fixés par le cahier des charges et la conduite du projet ;
  \item le \textbf{chapitre~\ref{chap:etat}} situe le travail par rapport aux méthodes classiques
        d'interpolation, aux modèles numériques de l'atmosphère et aux architectures d'apprentissage
        profond, dont celles que nous avons essayées puis abandonnées ;
  \item le \textbf{chapitre~\ref{chap:donnees}} décrit les données : le jeu synthétique ANEO, puis
        le domaine réel du Grand Tunis (relief, occupation du sol, ERA5, stations, radiosondages,
        prévisions GFS) ;
  \item le \textbf{chapitre~\ref{chap:modele}} détaille l'architecture PC-FNO3D, la perte physique,
        l'entraînement et la quantification de l'incertitude ;
  \item le \textbf{chapitre~\ref{chap:resultats}} rapporte les résultats, critère par critère, y
        compris les résultats négatifs et les décisions qui en ont découlé ;
  \item le \textbf{chapitre~\ref{chap:realisation}} présente la réalisation logicielle : paquet
        Python, API, serveur MCP, mise en service et application ATOM Wind ;
  \item la conclusion dresse le bilan et les perspectives.
\end{itemize}

Un principe a guidé tout le projet, celui de la preuve de concept de Nuthora Edu : \emph{un
résultat n'est retenu que s'il est mesuré face à des critères fixés à l'avance}. Cette exigence est
l'un des apports du travail : elle a permis d'écarter des pistes séduisantes mais fragiles
--- correction post-modèle, affinage sur les stations --- et de fonder chaque décision, y compris
les résultats négatifs, qui orientent directement la suite du projet.
